% hip-measured-boot.tex
% Hanzo Improvement Proposal — Measured Boot for the Hanzo microVM.
% Number assigned when this lands in HIPs/.
\documentclass[11pt,a4paper]{article}

\usepackage[utf8]{inputenc}
\usepackage[T1]{fontenc}
\usepackage[margin=1in]{geometry}
\usepackage{amsmath,amssymb,amsthm}
\usepackage{booktabs}
\usepackage{longtable}
\usepackage{array}
\usepackage{enumitem}
\usepackage{xcolor}
\usepackage{microtype}
\usepackage[colorlinks=true,linkcolor=black,citecolor=black,urlcolor=blue]{hyperref}

\newtheorem{definition}{Definition}[section]
\newtheorem{claim}[definition]{Claim}
\newtheorem{remark}[definition]{Remark}

% Long identifiers are the norm here, so let the typewriter font break.
\newcommand{\code}[1]{\texttt{\hyphenchar\font=`\-\relax #1}}
\sloppy
\emergencystretch=3em
\newcommand{\sha}[1]{\mathrm{SHA\text{-}#1}}
\newcommand{\zero}{\mathbf{0}}

\setlist[itemize]{itemsep=2pt,topsep=4pt}
\setlist[description]{itemsep=3pt,topsep=4pt}

\title{\textbf{Measured Boot: What a Machine Is, Stated as a Number}\\[6pt]
\large A Hanzo Improvement Proposal}

\author{Hanzo AI}
\date{6 September 2026}

\begin{document}
\maketitle

\begin{center}
\begin{tabular}{@{}ll@{}}
\toprule
Number      & assigned on merge into \code{HIPs/} \\
Type        & Standards Track \\
Category    & Security \\
Status      & Draft \\
Requires    & HIP-0005 (post-quantum security), HIP-0027 (KMS) \\
Referenced by & \code{hip-node-identity-wallet.tex}, \code{hip-confidential-tier.tex} \\
\bottomrule
\end{tabular}
\end{center}

\begin{abstract}
\noindent
A Hanzo microVM is started by a launcher on a host, boots a Linux kernel from a
root image, and runs a workload. This document specifies how that launch and
that workload are reduced to two numbers, how the numbers are recorded, how a
second party recomputes them from the same inputs, and how they are carried into
a hardware attestation report on platforms that produce one.

The construction is an ordered log of events folded into an extend-only SHA-384
register. It is chosen to match, exactly, the extend rule Intel TDX applies to a
runtime measurement register, so a register computed by the launcher is the value
an RTMR would hold had it started at zero and received the same events in the same
order. SHA-384 is the digest both TDX registers and the AMD SEV-SNP launch
measurement use.

Three properties are load-bearing and each is a consequence of the encoding rather
than of an implementation choice. An event's name is inside its own digest, so an
event cannot be replayed under a different name. The register commits to order, so
a prefix of a log and the whole log are distinct values. Only content is hashed ---
the filesystem path an image was read from is recorded for a reader and never
enters a digest --- so the same kernel and the same root image measure identically
on every machine that holds them, and a measurement is reproducible by any party
with the assets rather than being a fingerprint of one host's filesystem.

The document is explicit about what a measurement computed on a host does not
establish. It is a statement about what was launched, not evidence against the
host that launched it: a host able to substitute the kernel is also able to
substitute the program that hashes it. It becomes evidence only when bound to a
report signed by hardware the host does not control. Section~\ref{sec:hardware}
specifies that binding, states what hardware is required, and states precisely
which parts are implemented and which are not. None of the three machines the
implementation was built and measured on has the hardware, so the report path is
implemented and unexercised; Section~\ref{sec:status} says so, and says what the
tests do check in its place.
\end{abstract}

\tableofcontents

% =====================================================================
\section{Scope}
\label{sec:scope}

This document specifies the measurement of a Hanzo microVM: the event encoding,
the register fold, the division into a launch log and a workload log, the
document these are recorded in, the 64-byte value carried into a hardware report,
and the request by which a guest obtains that report.

It does not specify what a verifier does with a report --- the certificate chains,
the freshness terms, and the session binding are the subject of
\code{hip-confidential-tier.tex}, and identity is the subject of
\code{hip-node-identity-wallet.tex}. Where those documents and this one touch the
same object, this one defers: \S\ref{sec:field} states the relationship rather
than restating either.

Throughout, \emph{observed} marks a statement about code that was read or
behaviour that was measured, with the location or the number given. \emph{Proposed}
marks a change this document asks for. Statements with neither marker are
definitions or consequences of the two.

The key words MUST, MUST NOT, SHOULD, SHOULD NOT and MAY are to be interpreted as
in RFC 2119.

% =====================================================================
\section{Terms}
\label{sec:terms}

\begin{description}
\item[Launcher] The process that configures a hypervisor and starts a microVM.
  \emph{Observed}: the \code{hanzo-vm} binary, \code{hanzoai/vm},
  \code{crates/vm-cli}.

\item[Guest] The Linux system the launcher starts. Its first process is
  \code{vm-guest} (\code{crates/vm-guest}), reached from the host over vsock.

\item[Event] A named piece of content that was loaded. Written
  $e = (\mathit{name}, \mathit{source}, d)$ where $d$ is the SHA-384 of the
  content. \emph{source} records where the content came from or is going and is
  never hashed.

\item[Log] A finite ordered sequence of events.

\item[Register] The fold of a log, defined in \S\ref{sec:register}. Written $L$
  for the launch register and $W$ for the workload register.

\item[Launch] What the hypervisor was handed: kernel, kernel command line,
  initramfs where one is used, source root image, and shape. The launch log is
  the log of those events, in that order.

\item[Workload] What the started system was told to run. The workload log is
  extended by whatever drives the machine, after the guest is up.

\item[Shape] The machine the guest wakes up on, as the text
  \code{cpus=$N$ memory=$M$MB disk=$D$MB}.

\item[Source root image] The base image or checkpoint the per-instance working
  copy is cloned from. The working copy is written to during the run; the source
  is the object two parties can compare.

\item[Bind] The 64 bytes a hardware report carries as its caller-supplied field
  (REPORT\_DATA on SEV-SNP, the report data field on TDX). \S\ref{sec:bind}.

\item[Platform] A confidential-computing implementation that will sign a report
  for a guest it protects: AMD SEV-SNP or Intel TDX.
\end{description}

% =====================================================================
\section{The register}
\label{sec:register}

\subsection{Definition}

\begin{definition}[Event digest]
For an event $e = (\mathit{name}, \mathit{source}, d)$,
\[
  \varepsilon(e) \;=\; \sha{384}\bigl(\mathit{name} \,\|\, \mathtt{0x00} \,\|\, d\bigr),
\]
where $\mathit{name}$ is its UTF-8 bytes and $d$ is 48 bytes.
\end{definition}

\begin{definition}[Register]
For a log $\langle e_1, \dots, e_n \rangle$,
\[
  R_0 = \zero_{48}, \qquad
  R_i = \sha{384}\bigl(R_{i-1} \,\|\, \varepsilon(e_i)\bigr), \qquad
  R = R_n .
\]
\end{definition}

The recurrence is the extend rule TDX applies to an RTMR, so $R$ is the value an
RTMR would hold had it started at zero and been extended, in order, with each
$\varepsilon(e_i)$. \emph{Observed}: \code{crates/vm-measure/src/lib.rs},
\code{Log::register}; the rule is pinned by a test that computes one step by hand
rather than by calling the implementation.

A register is derived from its log on every read and never stored as an
independent field. \emph{Observed}: \code{Log} serializes as its value together
with its events, and deserializing recomputes the fold and refuses a document
whose stated value does not follow from its events.

\subsection{Three properties}

\begin{claim}[Names are not interchangeable]
\label{cl:name}
If $e$ and $e'$ have equal content digests and different names, then
$\varepsilon(e) \ne \varepsilon(e')$, except with the probability of a SHA-384
collision.
\end{claim}
\noindent The name is inside the preimage, and the $\mathtt{0x00}$ separator makes
the concatenation injective over a name set containing no NUL byte. Every name
this specification defines (\S\ref{sec:logs}) is a lowercase ASCII word. An
implementation MUST NOT admit a name containing $\mathtt{0x00}$.

\begin{claim}[Order is committed to]
\label{cl:order}
Two logs over the same multiset of events in different orders have different
registers, except with the probability of a SHA-384 collision.
\end{claim}

\begin{claim}[A prefix is not the whole]
\label{cl:prefix}
Extending a log always changes its register, except with the probability of a
SHA-384 collision.
\end{claim}
\noindent Claims~\ref{cl:order} and~\ref{cl:prefix} follow from $R_i$ being a
function of $R_{i-1}$; both are covered by tests.

\subsection{What is not hashed}
\label{sec:notpath}

An event records a \emph{source} for a reader --- the filesystem path an image was
read from, the path a manifest is written to, or the text itself when the content
is text --- and the source does not enter $\varepsilon$. Two consequences:

\begin{itemize}
\item The same kernel and the same root image measure identically on every machine
  that holds them, whatever they are called there. A measurement is therefore
  reproducible by any party holding the assets.
\item A file event's digest is an ordinary SHA-384 of the file's bytes, with no
  framing. \emph{Observed}: \code{shasum -a 384 \textasciitilde/.hanzo/vm/Image}
  prints the \code{kernel} event's digest, byte for byte.
\end{itemize}

The second is a deliberate constraint on the design and not a coincidence. A
measurement a second party cannot recompute with tools they already have is a
measurement they have to take on trust.

% =====================================================================
\section{The two logs}
\label{sec:logs}

A measurement is two logs, not one. They are extended by different programs at
different times --- the launcher, and then whatever drives the started machine ---
and on TDX they map onto separate runtime measurement registers. Folding them
into one chain would lose which half a verifier is entitled to predict.

\subsection{The launch log}

The launch log MUST contain exactly these events, in this order:

\begin{center}
\begin{tabular}{@{}lll@{}}
\toprule
Name & Content & Source recorded \\
\midrule
\code{kernel}  & the kernel image's bytes            & its path \\
\code{cmdline} & the kernel command line, UTF-8      & the same text \\
\code{initrd}  & the initramfs bytes \emph{(present only when one is used)} & its path \\
\code{root}    & the source root image's bytes       & its path \\
\code{shape}   & the shape text                      & the same text \\
\bottomrule
\end{tabular}
\end{center}

\emph{Observed}: \code{crates/vm-cli/src/measure.rs}, \code{launch}. Three
decisions in that list are worth stating explicitly.

The \emph{command line} is measured because it is part of what the guest is: the
same kernel told \code{root=/dev/vdb} is a different machine. It is produced by
one function (\code{vm::command\_line}) which is also the only thing that tells
the hypervisor, so the measured string and the handed string cannot differ.

The \emph{root} event measures the SOURCE image, never the per-instance working
copy. The copy is a clone of the source extended with zeros to the requested disk
size and is written to throughout the run; nothing about it is comparable between
two machines. The source is.

The \emph{shape} is measured because a platform that signs for a launch signs for
its vCPU count and its memory. A verifier that checked only the software would
accept the same kernel on a machine it never agreed to.

An \code{initrd} event is present only when an initramfs is actually used. The
alternative --- a zero digest when there is none --- would make ``no initramfs''
and ``an initramfs whose content hashes to zero'' the same measurement, and by
Claim~\ref{cl:prefix} the absence already moves the register.

\subsection{The workload log}

The workload log is extended by whatever drives the machine. This document fixes
its contents for exactly one driver, \code{hanzo up} (\code{hanzoai/cli},
\code{commands/up.rs}), which starts a k3s cluster and deploys the Hanzo cloud
into it:

\begin{center}
\begin{tabular}{@{}lll@{}}
\toprule
Name & Content & Source recorded \\
\midrule
\code{image}    & the pinned image reference, UTF-8 & the same text \\
\code{manifest} & the manifest text applied         & its path in the guest \\
\bottomrule
\end{tabular}
\end{center}

\code{image} is a digest reference (\code{registry/repository@sha256:\dots}),
resolved from whatever the caller named before the manifest is written. What the
manifest names, what the container runtime verifies on pull, and what the register
covers are then the same bytes by construction. \code{manifest} is the exact text
the cluster applies, because the same image deployed under a different Deployment
is a different cluster.

\begin{remark}[Per-instance state is excluded, on purpose]
The cluster's at-rest key --- 32 bytes from the operating system's generator, fresh
per cluster --- is written to the guest in a \emph{separate} file which is not
measured. A register is a statement about software identity; per-instance state in
it would make two clusters running identical software measure differently, which
defeats the reproducibility of \S\ref{sec:notpath}. \emph{Observed}: the two files
are \code{cloud-key.yaml} and \code{cloud.yaml}, and only the second is measured.
\end{remark}

An implementation extending the workload log for a different driver MUST NOT
include per-instance secrets, timestamps, or host-specific paths in any event's
content.

% =====================================================================
\section{The document}
\label{sec:document}

A measurement is recorded as one JSON document:

\begin{quote}\small
\begin{verbatim}
{ "algorithm": "sha384",
  "launch":   { "value": <hex 48>, "events": [ {name, source, digest}, ... ] },
  "workload": { "value": <hex 48>, "events": [ ... ] },
  "bind":     <hex 64>,
  "hardware": { "platform": "none" | "sev-snp" | "tdx",
                "report": <hex, present only with a platform> } }
\end{verbatim}
\end{quote}

Both registers and the bind are written for a reader who wants the numbers
without folding, and are \emph{verified} on the way back in: reading a document
recomputes each register from its events and recomputes the bind from the
registers, and refuses the document if either disagrees. A reader therefore never
has to decide whether to believe a stated value. \emph{Observed}:
\code{Measurement::from\_json}; a test edits a register and a bind in a serialized
document and asserts both are refused.

\code{hardware} is an observation carried alongside the chain and never folded
into it. \S\ref{sec:hardware}.

% =====================================================================
\section{The bind}
\label{sec:bind}

\begin{definition}[Bind]
\[
  \beta \;=\; \sha{512}\bigl(L \,\|\, W\bigr),
\]
64 bytes, over the two 48-byte registers.
\end{definition}

SHA-512 because the caller-supplied field in both report formats is exactly 64
bytes. A construction producing fewer would leave the remainder to be padded by
something, and padding nobody chose is a preimage nobody can reproduce.

\subsection{Who owns the caller field}
\label{sec:field}

A hardware report has one caller-supplied field, and more than one part of this
system would like to put something in it. \code{hip-confidential-tier.tex}
specifies its contents for a served-inference session (its requirement B1a),
fixing each term to a width and including the session key-exchange key, the
receipt verifying key, a session identifier and a verifier nonce alongside $L$ and
$W$. That specification governs wherever a session exists. This document does not
propose a second.

What this document specifies is narrower and, by construction, compatible: the
launcher's request interface takes 64 \emph{opaque} bytes and asks the platform
for a report over exactly those. The launcher has no opinion about their
derivation. $\beta$ is what a caller with nothing but a boot to describe should
send; a caller with a session sends B1a's value through the same interface.

\emph{Observed}: the vsock request payload is 64 raw bytes with no structure, and
the guest hands them to the platform unmodified
(\code{crates/vm-guest/src/main.rs}, \code{ATTEST\_REQ}).

\begin{remark}[What $\beta$ alone establishes]
$\beta$ is a function of the release. Two machines running the same launch and the
same workload produce the same $\beta$, and a report over it carries no freshness
term. \code{hip-node-identity-wallet.tex} states both consequences formally (its
Claims on transferability and staleness) and neither is disputed here: a report
over $\beta$ establishes that \emph{a} platform of the stated kind launched
\emph{this} image at \emph{some} time. That is worth having and it is not
identity, and it is why the session terms in B1a exist.
\end{remark}

% =====================================================================
\section{The hardware request}
\label{sec:hardware}

\subsection{Where the request is made, and why there}

Both platforms expose the same primitive: a guest asks its firmware for a report
over 64 caller-supplied bytes, and the report also carries the platform's own
measurement of how that guest was started. The request is made through a guest
device --- \code{/dev/sev-guest} or \code{/dev/tdx\_guest} --- which exists only
inside a guest the platform protects.

The request MUST therefore be made from inside the guest. A request made on the
host answers a different question: whether the machine running the launcher is
itself somebody's confidential guest. A report obtained that way, over our bind,
says nothing about the VM we started.

\emph{Observed}: \code{vm\_measure::attest::status} is called by \code{vm-guest},
reached over vsock. The launcher's \code{hanzo-vm measure}, which states what a
boot \emph{would} be and starts nothing, records \code{platform: none} and asks no
device --- there is no guest to ask.

\subsection{The wire}

\begin{center}
\begin{tabular}{@{}lll@{}}
\toprule
Frame & Type byte & Payload \\
\midrule
\code{ATTEST\_REQ}  & \code{0x60} & the 64 bytes, raw \\
\code{ATTEST\_RESP} & \code{0x61} & JSON \code{\{platform, report?\}} \\
\bottomrule
\end{tabular}
\end{center}

A guest receiving an \code{ATTEST\_REQ} whose payload is not exactly 64 bytes MUST
answer with an error frame rather than pad or truncate: the payload is the field
the report is taken over, and either adjustment produces a report over bytes
nobody chose.

A guest that predates this frame drops it in silence (its dispatch ignores unknown
types). A host MUST therefore bound the read. \emph{Observed}: five seconds,
\code{Sandbox::attest}.

\subsection{The two ABIs}

\emph{Observed}, \code{crates/vm-measure/src/attest.rs}. The request structures and
ioctl numbers are those of \code{linux/sev-guest.h} and \code{linux/tdx-guest.h}:

\begin{center}
\begin{tabular}{@{}llll@{}}
\toprule
Platform & Device & ioctl & Report \\
\midrule
SEV-SNP & \code{/dev/sev-guest}  & \code{SNP\_GET\_REPORT}, \code{0xC0205300}     & 1184 B \\
TDX     & \code{/dev/tdx\_guest} & \code{TDX\_CMD\_GET\_REPORT0}, \code{0xC4405401} & 1024 B \\
\bottomrule
\end{tabular}
\end{center}

Both numbers are computed in the source by the \code{\_IOWR} encoding rather than
copied as literals, so a structure that drifts from the kernel's would produce a
number no driver answers, and the drift is caught rather than sent. The report is
recorded exactly as the platform produced it; parsing belongs to the verifier.

Detection is device presence and nothing else --- the drivers bind only inside a
guest the platform actually protects. A device that will not answer records
\code{none} rather than asserting a platform that produced nothing.

\subsection{What hardware is required}
\label{sec:hw}

The path requires, together:

\begin{enumerate}[label=(\roman*)]
\item A CPU with the feature: AMD EPYC, Milan (3rd generation) or later, with
  SEV-SNP enabled in firmware; or an Intel Xeon with TDX enabled.
\item A VMM that starts guests of that kind --- for SEV-SNP, cloud-hypervisor
  driving an IGVM image, or QEMU with SNP support; for TDX, a TDX-aware VMM.
\item A guest kernel carrying the corresponding driver, so the device node exists.
\end{enumerate}

None of the three machines the implementation was built and measured on satisfies
(i). Apple Silicon has no third-party equivalent; NVIDIA's Grace-Blackwell
workstation parts have none exposed to a VMM of ours; the x86 workstation is not
an SEV-SNP part. The launcher's three backends --- Virtualization.framework, the
arm64 KVM backend, and cloud-hypervisor as we drive it --- all start ordinary
guests, so (ii) is unmet as well.

\emph{Proposed.} Reaching an exercised path requires exactly two things and
neither is in this document's tree: an SEV-SNP-capable host in the fleet, and the
launcher's cloud-hypervisor backend configured to start an SNP guest from an IGVM
image. Everything above it --- the request, the wire, the guest handler, the
document field --- is in place and would begin returning reports without a change.

% =====================================================================
\section{Reproducibility and cost}
\label{sec:cost}

\subsection{Recomputing a measurement}

A second party holding the same kernel, root image, command line and shape
recomputes both registers with a SHA-384 implementation and nothing else. The
file digests are checkable with \code{shasum}; the fold is four lines.

\subsection{The digest cache}

A root image is gigabytes. \emph{Observed}: hashing a 4~GB root image takes
19.3~s; answering from a remembered digest takes 9~ms. A launcher that re-read the
image on every boot would cost more in measurement than in boot.

Digests are therefore remembered against the identity the filesystem reports for
each file --- device, inode, length, modification time --- and recomputed the
moment any of those move. The cache is warmed where a slow path is already being
paid: after an image is downloaded, and after a checkpoint is saved.

The cache carries no authority and is not part of this specification's security
argument. It lives on the same disk as the images it describes, so a party able to
rewrite one is able to rewrite the other; it saves time and proves nothing. An
implementation MUST provide a way to bypass it (\emph{observed}:
\code{hanzo-vm measure --recompute}), and a verifier reproducing a measurement
MUST NOT use one.

% =====================================================================
\section{Security considerations}
\label{sec:security}

\subsection{What a host-computed measurement is}

A measurement computed by a launcher is a statement about what the launcher
loaded. A host able to substitute the kernel is also able to substitute the
program that hashes it, so the measurement is not evidence \emph{against} the host
that produced it. This is stated once in the source, at the top of the crate, and
is repeated here because it is the boundary the rest of the design is arranged
around.

What it is good for, on its own:

\begin{itemize}
\item \emph{Reproducibility.} Two parties can agree on what a release is, and
  disagree with a machine that claims to be running it.
\item \emph{Change detection.} A register that moves when nothing was meant to
  change is a fact about the machine, whoever computed it.
\item \emph{A predicate a verifier can state in advance.} An expected register can
  be published before a machine is asked for one.
\end{itemize}

What it is not: an authenticated statement. That requires \S\ref{sec:hardware},
and on hardware that does not exist in our fleet today.

\subsection{The measurement and the workload}

The workload register commits to what the driver \emph{deployed}, not to what the
cluster subsequently ran. A cluster whose control plane is reachable can be given
a second workload afterwards, and no register moves. The register bounds the
cluster's initial state; continuous integrity of a running cluster is a different
mechanism and is out of scope.

\subsection{The digest}

SHA-384 is chosen for interoperability with the two platform formats
(\S\ref{sec:register}) and not because 384 bits is required. It is truncated
SHA-512 and therefore not vulnerable to length extension, which matters here
because $\varepsilon$ and the fold both hash concatenations. The bind uses
SHA-512 for the width of the field it fills. Neither choice is post-quantum; a
measurement is a commitment rather than a signature, so the relevant property is
collision resistance, and SHA-384's is not reduced below any threshold this system
depends on by known quantum algorithms.

% =====================================================================
\section{Implementation status}
\label{sec:status}

\emph{Observed}, \code{hanzoai/vm} v2.0.1 and \code{hanzoai/cli} v8.5.158.

In \code{hanzoai/vm} (paths relative to \code{crates/}):

\begin{center}
\begin{tabular}{@{}lll@{}}
\toprule
Part & Where & State \\
\midrule
Encoding, fold, document    & \code{vm-measure/src/lib.rs}    & implemented, tested \\
Digest cache                & \code{vm-measure/src/lib.rs}    & implemented, tested \\
SEV-SNP / TDX ioctls        & \code{vm-measure/src/attest.rs} & implemented, unexercised \\
Launch log                  & \code{vm-cli/src/measure.rs}    & implemented, tested \\
\code{hanzo-vm measure}     & \code{vm-cli/src/main.rs}       & implemented \\
Launch report over the wire & \code{vm-cli/src/stdio.rs}      & implemented \\
Report request, host side   & \code{vm/src/sandbox.rs}        & implemented \\
Report request, guest side  & \code{vm-guest/src/main.rs}     & implemented \\
\bottomrule
\end{tabular}
\end{center}

\noindent In \code{hanzoai/cli}, all of it in \code{src/commands/up.rs}: the
workload log, the document filed at \code{\textasciitilde/.hanzo/up/measure.json},
and \code{hanzo up --attest}. Implemented and tested.

``Unexercised'' is exact: no machine available to the implementation can make the
ioctl succeed (\S\ref{sec:hw}). What is checked in its place is the part a machine
without the hardware can check --- the three request structures are asserted to be
the sizes the kernel headers define (96, 32 and 1088 bytes), and both ioctl numbers
are asserted to equal the documented values. An ioctl number encodes the size of
the structure it carries, so a structure that drifted would send a number no driver
answers, and that is the failure these tests exist to catch.

% =====================================================================
\section{Conformance}
\label{sec:conformance}

An implementation conforms to this document if:

\begin{enumerate}[label=C\arabic*.,leftmargin=2.5em]
\item It computes $\varepsilon$ and the fold as defined in \S\ref{sec:register},
  from a 48-byte zero register.
\item Its launch log carries the events of \S\ref{sec:logs} with those names, in
  that order, and includes \code{initrd} exactly when an initramfs is used.
\item It measures the source root image, not a per-instance copy.
\item No event's content includes a filesystem path, a per-instance secret, or a
  timestamp.
\item It derives every register from its log on read and refuses a document whose
  stated register or bind does not follow.
\item It computes $\beta$ as in \S\ref{sec:bind} and requests a report over 64
  bytes it does not interpret.
\item It makes the report request inside the guest, and records \code{none} where
  no platform device is present rather than asserting one.
\item It provides a way to compute a measurement without a digest cache.
\end{enumerate}

% =====================================================================
\section{Unsolved}
\label{sec:unsolved}

\begin{enumerate}[label=U\arabic*.,leftmargin=2.5em]
\item \textbf{No hardware.} \S\ref{sec:hw}. Until the fleet holds an SEV-SNP or
  TDX host and the launcher can start such a guest, every document this system
  produces records \code{platform: none}, and the measurement is a reproducibility
  mechanism rather than an evidence one.

\item \textbf{The launcher is not measured.} The launch log covers what the
  hypervisor was handed, not the hypervisor. On a confidential VM that is
  correct --- the platform's own launch measurement covers the VMM-supplied state,
  so measuring the launcher from inside would duplicate a value the hardware
  already signs. Where there is no platform, nothing covers it, and a host that
  swaps its launcher binary produces the same register.

\item \textbf{Nothing binds a measurement to a machine.} $\beta$ is a function of
  the release; two hosts running the same release produce the same value. The
  binding of a measurement to an identity is
  \code{hip-node-identity-wallet.tex}'s subject and is not attempted here.

\item \textbf{The workload register is a commitment to a first state.}
  \S\ref{sec:security}. A cluster that is subsequently changed does not move a
  register.
\end{enumerate}

\end{document}
